% Paquete de remisión CEI UDIT — anexos del comunicado (blend)
% Build: make remision  (XeLaTeX — CJK for 道)
\documentclass[11pt,a4paper]{scrartcl}

\usepackage[spanish]{babel}
\usepackage{amssymb}
\usepackage{microtype}
\usepackage{geometry}
\usepackage{booktabs}
\usepackage{setspace}
\IfFileExists{enumitem.sty}{%
  \usepackage{enumitem}%
  \setlist[itemize]{leftmargin=1.4em, itemsep=0.12em, topsep=0.15em}%
  \setlist[enumerate]{leftmargin=1.6em, itemsep=0.12em, topsep=0.15em}%
}{%
  \setlength{\leftmargini}{1.4em}%
}

% Shared chrome for CEI / student ethics PDFs — UDIT · ECSIT
% Built with XeLaTeX (CJK for 道). Input from documents in this directory.
\usepackage{graphicx}
\usepackage{fontspec}
\usepackage{hyperref}
\usepackage{xcolor}
\usepackage{tabularx}
\usepackage{fancyhdr}
\usepackage{calc}

% Single CJK glyph (道) via fontspec — no xeCJK dependency
\IfFontExistsTF{Songti SC}{%
  \newfontfamily\TTODdao{Songti SC}%
}{%
  \IfFontExistsTF{PingFang SC}{%
    \newfontfamily\TTODdao{PingFang SC}%
  }{%
    \IfFontExistsTF{STSong}{%
      \newfontfamily\TTODdao{STSong}%
    }{%
      \newcommand{\TTODdao}{}%
    }%
  }%
}

\definecolor{uditnavy}{HTML}{0B1F33}
\definecolor{uditink}{HTML}{1A1A1A}
\definecolor{rulegray}{HTML}{CCCCCC}

\newcommand{\assetspath}{../assets/}
\newcommand{\EcsitUrl}{https://www.udit.es/lineas-de-investigacion/grupo-de-investigacion-innovacion-y-tecnologia-desde-y-para-la-educacion-la-cultura-y-la-sociedad/}
\newcommand{\EcsitLink}{\href{\EcsitUrl}{udit.es $\rightarrow$ grupo ECSIT}}
\newcommand{\TTODfull}{{\TTODdao 道} The Tao of Development (TTOD)}

\newcommand{\LogoRow}{%
  \noindent
  \begin{tabularx}{\textwidth}{@{}lXr@{}}
    \includegraphics[height=0.95cm]{\assetspath udit-logo-print.png}
    & &
    \includegraphics[height=1.45cm]{\assetspath ecsit-logo-print.png}
  \end{tabularx}%
  \par\vspace{0.55em}
  \noindent{\color{rulegray}\hrule height 0.35pt}%
  \par\vspace{0.65em}
}

\newcommand{\LegalEntities}{%
  \noindent{\scriptsize\color{uditink}
  UDIT ESTUDIOS SUPERIORES INTERNACIONALES, S.L.\ · NIF B84014265 ·
  UDIT Escuela de Formación Profesional, S.L.\ · NIF B09963117.}%
}

% Role (left) | Name + email only (right) — no role text repeated
\newcommand{\AuthoritiesBlock}{%
  \noindent{\sffamily\bfseries\small Autoridades de investigación / Research authorities}\\[0.25em]
  \noindent
  \begin{tabularx}{\textwidth}{@{}>{\raggedright\arraybackslash}p{4.6cm}X@{}}
    \textbf{Coordinación grados tecnológicos} &
      Sandra Garrido ---
      \href{mailto:sandra.garrido@udit.es}{sandra.garrido@udit.es}\\[0.2em]
    \textbf{Dirección de grados} &
      Fernando Blázquez Piñeiro ---
      \href{mailto:fernando.blazquez@udit.es}{fernando.blazquez@udit.es}\\[0.2em]
    \textbf{IP del grupo ECSIT} &
      Rafael Conde Melguizo, PhD ---
      \href{mailto:rafael.conde@udit.es}{rafael.conde@udit.es}\\[0.2em]
    \textbf{CEI UDIT} &
      \href{mailto:comite.etica.investigacion@udit.es}{comite.etica.investigacion@udit.es}
  \end{tabularx}%
  \par\vspace{0.35em}%
  \LegalEntities
}

\newcommand{\StudyMetaPI}{%
  \noindent
  \begin{tabularx}{\textwidth}{@{}>{\raggedright\arraybackslash}p{4.6cm}X@{}}
    \textbf{Investigador principal} &
      Rubén Vega Balbás, PhD ---
      \href{mailto:ruben.vega@udit.es}{ruben.vega@udit.es}\\[0.15em]
    \textbf{Afiliación académica} &
      Grado Oficial en Desarrollo Full-Stack
      (\emph{La evolución lógica de la ingeniería informática hacia una
      titulación enfocada en el mundo real.})\\[0.15em]
    \textbf{Proyecto / producto} &
      \TTODfull\ · plataforma Oráculo · Cohorte Front-End~II\\[0.15em]
    \textbf{Grupo de investigación} &
      ECSIT · \EcsitLink\\[0.15em]
    \textbf{Comité} &
      \href{mailto:comite.etica.investigacion@udit.es}{comite.etica.investigacion@udit.es}\\[0.15em]
    \textbf{Código de estudio} &
      TTOD-FEII-2026-27\\[0.15em]
    \textbf{Versión del paquete} &
      2026-09-26 · cohorte pedagógica $n=8$
      (uso investigador: solo con consentimiento, $n\leq 8$)
  \end{tabularx}%
}

% Full Chicago (author–date) methods backbone — from docs/public/research/methodology.md
\newcommand{\ChicagoMethodsRefs}{%
\begin{itemize}
  \item Brown, Neil C.~C., and Mark Guzdial. 2024.
        ``Confidence vs Insight: Big and Rich Data in Computing Education Research.''
        In \emph{Proceedings of the 55th ACM Technical Symposium on Computer Science Education V.~1},
        158--64. \url{https://doi.org/10.1145/3626252.3630813}.
  \item Fischer, Björn, Berit Barthelmes, Sven Eric Panitz, Eva-Maria Iwer, and Ralf Dörner. 2024.
        ``Seeking Consent for Programming Process Data Collection with Trustee-Based Encryption.''
        In \emph{Proceedings of the 2024 ACM Conference on International Computing Education Research V.1},
        131--42. \url{https://doi.org/10.1145/3632620.3671125}.
  \item Heckman, Sarah, Jeffrey C.~Carver, Mark Sherriff, and Ahmed Al-Zubidy. 2022.
        ``A Systematic Literature Review of Empiricism and Norms of Reporting in Computing Education Research Literature.''
        \emph{ACM Transactions on Computing Education} 22 (1): Article~3.
        \url{https://doi.org/10.1145/3470652}.
  \item McGill, Monica M., Sarah Heckman, Christos Chytas, et~al. 2023.
        ``Conducting Sound, Equity-Enabling Computing Education Research.''
        In \emph{Proceedings of the 2023 Working Group Reports on Innovation and Technology in Computer Science Education},
        30--56. \url{https://doi.org/10.1145/3623762.3633495}.
  \item Runeson, Per, and Martin Höst. 2009.
        ``Guidelines for Conducting and Reporting Case Study Research in Software Engineering.''
        \emph{Empirical Software Engineering} 14 (2): 131--64.
        \url{https://doi.org/10.1007/s10664-008-9102-8}.
  \item Runeson, Per, Martin Höst, Austen Rainer, and Björn Regnell. 2012.
        \emph{Case Study Research in Software Engineering: Guidelines and Examples}.
        Hoboken, NJ: Wiley. \url{https://doi.org/10.1002/9781118181034}.
  \item Wohlin, Claes, Per Runeson, Martin Höst, Magnus C.~Ohlsson, Björn Regnell, and Anders Wesslén. 2012.
        \emph{Experimentation in Software Engineering}. Berlin: Springer.
        \url{https://doi.org/10.1007/978-3-642-29044-2}.
\end{itemize}%
}


\geometry{margin=2.2cm, top=2.0cm, bottom=2.2cm}
\setkomafont{section}{\sffamily\bfseries\large}
\setkomafont{subsection}{\sffamily\bfseries\normalsize}
\setkomafont{disposition}{\sffamily}
\setlength{\parskip}{0.32em}
\setlength{\parindent}{0pt}
\setstretch{1.04}

\hypersetup{
  colorlinks=true,
  urlcolor=uditnavy,
  linkcolor=uditnavy,
  pdftitle={Paquete de remisión CEI UDIT — 道 The Tao of Development},
  pdfauthor={Rubén Vega Balbás, PhD · UDIT / ECSIT},
  pdfsubject={Anexos del comunicado del Comité de Ética}
}

\pagestyle{fancy}
\fancyhf{}
\renewcommand{\headrulewidth}{0pt}
\fancyfoot[L]{\scriptsize\color{uditnavy} Remisión CEI · \TTODfull}
\fancyfoot[C]{\scriptsize\color{uditink}\thepage}
\fancyfoot[R]{\scriptsize\color{uditink}v.\ 2026-09-26}

\newcommand{\AnnexTitle}[1]{%
  \clearpage
  \LogoRow
  \begin{center}{\Large\sffamily\bfseries #1}\end{center}
  \vspace{0.4em}
  \noindent{\color{rulegray}\hrule height 0.4pt}
  \vspace{0.5em}
}

\begin{document}

%======================================================================
% PORTADA
%======================================================================
\LogoRow

\begin{center}
  {\LARGE\sffamily\bfseries Paquete de remisión}\\[0.35em]
  {\Large\sffamily Comité de Ética en Investigación (CEI) --- UDIT}\\[0.8em]
  {\large\sffamily Aprendizaje haciendo --- grupo \TTODfull\ / Front-End~II}\\[0.3em]
  {\normalsize\color{uditnavy} Estudio de caso educativo / investigación basada en diseño}\\[0.6em]
  {\small Código TTOD-FEII-2026-27 · Fecha de presentación: 26 septiembre 2026}
\end{center}

\vspace{0.6em}
\StudyMetaPI

\vspace{0.55em}
\AuthoritiesBlock

\vspace{0.7em}
\noindent{\sffamily\bfseries\small Documentos de este PDF (requisitos del comunicado)}\\[0.3em]
\begin{enumerate}
  \item Carta / solicitud narrativa al CEI
  \item Protocolo de investigación completo
  \item Instrumentos de recolección de datos
  \item Plan de manejo de datos y minimización
  \item Autorizaciones institucionales (checklist + pendientes)
  \item Hoja de información y formulario de consentimiento (plantilla estudiante)
  \item Declaración de exclusión de la línea filosófica paralela
\end{enumerate}

\vspace{0.4em}
\noindent{\small\color{uditink}
\textbf{Remisión completa =} este PDF +
\textbf{Formulario de Solicitud de Evaluación Ética} (Word institucional relleno:
\texttt{FORMULARIO-SOLICITUD-CEI-UDIT-RELLENO.docx}) +
adjuntos firmados fuera de git (autorización departamental; consentimientos firmados
según vía del CEI).}

\vspace{0.5em}
\LegalEntities

%======================================================================
% 1. CARTA
%======================================================================
\AnnexTitle{1.\ Carta de solicitud al CEI}

Se solicita evaluación ética de un \textbf{estudio de caso educativo / investigación
basada en diseño}, acotado a Front-End~II (curso 2026--27). El alumnado participa
como desarrolladores en un producto real (plataforma Oráculo de \TTODfull), con asistencia
de IA agentica y trabajo en grupo. La actividad docente fue autorizada
pedagógicamente con antelación. La participación en investigación (uso secundario
seudonimizado de artefactos de proceso) es \textbf{voluntaria} e independiente de
la calificación. La autorización de imagen/comunicación (ítem~G) es opcional y separable.

\textbf{No se pretende} demostrar causalmente que la IA <<mejora el aprendizaje>>.
Se pretende describir \textbf{cómo se aprende haciendo} y qué principios de diseño
pedagógico se desprenden.

\textbf{Pregunta principal:} ¿Puede la participación como desarrollador/a en una
aplicación real, con asistencia agentica y en grupo, constituir una actividad de
aprendizaje con resultados pedagógicos observables?

\textbf{Decisión solicitada:}
\begin{enumerate}
  \item Evaluar y, en su caso, autorizar el protocolo descrito.
  \item Confirmar que el enfoque de minimización / datos personales no puestos
        <<en juego>> como objeto de interés es adecuado, o indicar condiciones.
  \item Indicar si se requiere dictamen del DPD en paralelo
        (\href{mailto:dpd@udit.es}{dpd@udit.es}).
\end{enumerate}

%======================================================================
% 2. PROTOCOLO
%======================================================================
\AnnexTitle{2.\ Protocolo de investigación completo}

\subsection*{Tipo de estudio}
Estudio de caso educativo / investigación basada en diseño (cualitativo dominante,
datos ricos de proceso). \textbf{No} es ensayo clínico, \textbf{no} es experimento
controlado, \textbf{no} pretende tamaños de efecto ni inferencia poblacional.

\subsection*{Antecedentes y justificación}
La asistencia cotidiana de IA generativa en el desarrollo de software obliga a
distinguir \emph{código producido} de \emph{aprendizaje demostrado}, en especial
en cursos post-introductorios. Este estudio observa un grupo real que construye
un producto auténtico con asistencia agentica, trabajo colaborativo y defensa oral.

Justificación metodológica: caso acotado con evidencia de proceso (Runeson y Höst
2009; Runeson et al.\ 2012), datos ricos frente a grandes (Brown y Guzdial 2024),
normas CER (Heckman et al.\ 2022; McGill et al.\ 2023), consentimiento para datos
de proceso de programación (Fischer et al.\ 2024). Wohlin et al.\ (2012) sitúa el
diseño como \textbf{estudio de caso}, no como experimento.

\subsection*{Objetivos}
\textbf{General:} comprender si --- y cómo --- la participación como desarrollador/a
en un producto real, asistido por IA agentica y en grupo, constituye
\emph{aprender haciendo} con resultados pedagógicos observables.

\textbf{Específicos:}
\begin{enumerate}
  \item Identificar trazas de proceso que hacen explicables y revisables las
        contribuciones asistidas por IA.
  \item Describir negociación de autoría, agencia y responsabilidad.
  \item Proponer principios de diseño situados (no marco causal validado).
  \item Documentar divergencias entre producto entregado y aprendizaje demostrado.
\end{enumerate}

\subsection*{Diseño y método}
\renewcommand{\arraystretch}{1.2}
\begin{tabularx}{\textwidth}{@{}lX@{}}
  \toprule
  Diseño & Caso único acotado / DBR cualitativo dominante \\
  Caso & Cohorte Front-End II 2026--27 · producto Oráculo de \TTODfull \\
  Unidad de análisis & Trayectorias de pareja/solo (PRs) en el sistema del grupo \\
  Muestreo & Censo del grupo; $n=8$ estudiantes-desarrolladores \\
  Análisis & Codificación temática; triangulación producto / registro / defensa \\
  Inicio del análisis & Tras dictamen ético + consentimiento + cierre de notas \\
  \bottomrule
\end{tabularx}

\subsection*{Participantes}
\begin{itemize}
  \item \textbf{Inclusión:} matriculados en Front-End~II en la entrega TTOD que
        otorgan consentimiento de uso investigador.
  \item \textbf{Exclusión del corpus:} no consentidos; menores/vulnerables sin
        salvaguarda adicional; artefactos de compañeros no consentidos.
  \item \textbf{n:} 8 (README del proyecto; excluido el PI). Corpus: $n_{\mathrm{consent}}\leq 8$.
  \item \textbf{Reclutamiento:} no externo; la actividad es la asignatura.
\end{itemize}

\subsection*{Procedimiento}
\begin{enumerate}
  \item Actividad docente ordinaria (desarrollo, PRs, declaración de IA, defensa).
  \item Información y consentimiento (docente + investigador).
  \item Tras dictamen y consentimientos: corpus \textbf{seudonimizado}.
  \item Análisis cualitativo / de diseño.
  \item Difusión agregada o extractos no identificativos (consentimiento de difusión).
\end{enumerate}
Instrumentos opcionales (encuesta, entrevista, grabación) \textbf{no} forman parte
de esta versión.

\subsection*{Aspectos éticos (síntesis)}
Voluntariedad; nota no condicionada; separación docente/investigador; minimización;
sin categorías especiales; retirada sin consecuencia académica; línea filosófica
paralela \textbf{excluida} (Anexo~7).

\subsection*{Cronograma orientativo}
Docencia curso 2026--27 · dictamen tras remisión · análisis tras notas · difusión
académica (JENUI~$\rightarrow$~ITiCSE) tras análisis.

\subsection*{Referencias principales (Chicago author--date)}
\noindent{\small Espina metodológica --- títulos completos; misma lista canónica que
\texttt{docs/public/research/methodology.md}.}
\ChicagoMethodsRefs

%======================================================================
% 3. INSTRUMENTOS
%======================================================================
\AnnexTitle{3.\ Instrumentos de recolección de datos}

\subsection*{Principio}
Este estudio \textbf{no introduce} cuestionarios, baterías psicométricas, diarios
de vida privada ni entrevistas como instrumentos por defecto. La recolección
investigadora es \textbf{uso secundario seudonimizado} de artefactos ya producidos
para la evaluación ordinaria, tras consentimiento informado (ítem~B del formulario).

\subsection*{Instrumentos activos (diseño mínimo)}
\renewcommand{\arraystretch}{1.25}
\begin{tabularx}{\textwidth}{@{}c>{\raggedright\arraybackslash}X>{\raggedright\arraybackslash}p{3.2cm}@{}}
  \toprule
  \textbf{ID} & \textbf{Instrumento / origen} & \textbf{Formato} \\
  \midrule
  I1 & Historial commits / PRs (remoto del curso) --- narrativas y revisiones seudonimizadas & Texto / metadatos \\
  I2 & Declaración de uso de IA (rúbrica FE~II) --- sin datos de vida privada & Texto \\
  I3 & Notas estructuradas de defensa oral --- \textbf{sin} audio/vídeo & Texto \\
  I4 & Extractos de código no identificativos --- solo con consentimiento de difusión & Texto / imagen \\
  \bottomrule
\end{tabularx}

\subsection*{Instrumentos no activados}
I5 cuestionario · I6 entrevista · I7 grupo focal · I8 grabación · I9 telemetría IDE /
keystrokes / prompts completos --- \textbf{N/A} en esta versión. Su activación
exigiría modificación al CEI y consentimiento ampliado.

\subsection*{Procedimiento investigador}
\begin{enumerate}
  \item El alumnado produce I1--I3 como parte de la docencia.
  \item Tras dictamen y consentimiento, se construye corpus seudonimizado (\texttt{P01\ldots Pn}).
  \item No se exportan calificaciones ni identificadores al corpus.
\end{enumerate}

%======================================================================
% 4. PLAN DE DATOS
%======================================================================
\AnnexTitle{4.\ Plan de manejo de datos y minimización}

\subsection*{Resumen}
Datos de investigación = trazas de proceso y textos pedagógicos seudonimizados.
\textbf{No} cuestionarios ni entrevistas. Notas e identidades académicas:
sistemas docentes, \textbf{fuera} del corpus.

\subsection*{Conjuntos}
\begin{itemize}
  \item Corpus investigador (I1--I4): seudónimos; sin nombre, email, DNI.
  \item Mapa seudónimo$\leftrightarrow$identidad: offline, cifrado; no en git.
  \item Consentimientos firmados: custodia privada.
  \item Categorías especiales (art.~9 RGPD): \textbf{no} se recogen.
\end{itemize}

\subsection*{Base jurídica (propuesta; confirmar con DPD)}
Docencia/evaluación: marco académico institucional.
Uso investigador secundario: \textbf{consentimiento} (art.~6.1.a RGPD / LOPDGDD).

\subsection*{Almacenamiento y seguridad}
Corpus: almacenamiento institucional / cifrado, acceso mínimo, sin nube pública
de datos de estudiantes. Mapa seudónimo separado. Prohibido subir consentimientos
firmados o mapas nominativos al repositorio git del proyecto.

\subsection*{Conservación (propuesta)}
Consentimientos: fin del estudio + plazo institucional ($\geq 5$~años o norma UDIT).
Corpus: hasta cierre de publicaciones del ciclo. Mapa seudónimo: destruir cuando
deje de ser necesario. Confirmación DPD:
\href{mailto:dpd@udit.es}{dpd@udit.es} ·
\href{mailto:proteccion.datos@udit.es}{proteccion.datos@udit.es}.

\subsection*{Eliminación}
Eliminación digital segura; destrucción de documentos físicos si existen en papel.

\subsection*{Postura de minimización (síntesis)}
El estudio está diseñado para que los datos personales del alumnado \textbf{no
estén en juego} como objeto de interés investigador: sin encuestas de vida privada,
sin categorías especiales, sin publicación de identificadores. Se reconoce que, bajo
RGPD, datos seudonimizados pueden seguir siendo personales; por ello el corpus
es confidencial y \textbf{no} se trata como dato anónimo abierto.

%======================================================================
% 5. AUTORIZACIONES
%======================================================================
\AnnexTitle{5.\ Autorizaciones institucionales}

\renewcommand{\arraystretch}{1.25}
\begin{tabularx}{\textwidth}{@{}c>{\raggedright\arraybackslash}X>{\raggedright\arraybackslash}p{3.4cm}@{}}
  \toprule
  \textbf{ID} & \textbf{Autorización} & \textbf{Estado} \\
  \midrule
  A1 & Encuadre pedagógico departamental de la actividad del grupo & Afirmada por el PI; \textbf{adjunto firmado pendiente} \\
  A2 & Dictamen del CEI (esta remisión) & En curso \\
  A3 & Consulta / visto bueno DPD & Pendiente / a criterio institucional \\
  A4 & Autorización adicional de centro/grado si el reglamento lo exige & Confirmar \\
  A5 & Convenios con terceros / cloud de datos de investigación & N/A \\
  \bottomrule
\end{tabularx}

\vspace{0.5em}
\noindent
La autorización \textbf{pedagógica} (A1) \textbf{no} equivale a la autorización
\textbf{ética de investigación} (A2). Ambas deben constar en el expediente.

\noindent\textbf{Adjuntos a aportar fuera de este PDF cuando existan:}
correo/acta departamental firmada; extracto de guía docente / encargo;
respuesta DPD si procede; consentimientos firmados según vía del CEI.

%======================================================================
% 6. CONSENTIMIENTO (plantilla)
%======================================================================
\AnnexTitle{6.\ Información y consentimiento informado (plantilla)}

\noindent{\small
Plantilla para participantes. La versión tipográfica estudiante (mismos logos)
también se distribuye como PDF separado
\texttt{HOJA-INFORMACION-Y-CONSENTIMIENTO-ES.pdf}. Firmados: \textbf{fuera de git}.}

\subsection*{Síntesis de la hoja de información}
\begin{itemize}
  \item Estudio: aprendizaje haciendo en grupo \TTODfull\ / Front-End~II.
  \item Docencia sigue igual con o sin consentimiento investigador ni autorización de imagen.
  \item Investigación = uso seudonimizado de artefactos de evaluación (PRs,
        declaración de IA, notas de defensa), tras dictamen y cierre de notas.
  \item Ítem~G (opcional): imagen en LinkedIn/medios/publicaciones/revistas internas UDIT/web público.
  \item Sin encuestas de vida privada, sin categorías especiales; grabación de
        defensa/entrevista N/A en el diseño mínimo (distinta de G).
  \item Voluntario; retirada sin efecto en la calificación.
\end{itemize}

\vspace{0.3em}
\AuthoritiesBlock

\subsection*{Ítems del formulario}
\renewcommand{\arraystretch}{1.3}
\begin{tabularx}{\textwidth}{@{}c>{\raggedright\arraybackslash}X c@{}}
  \toprule
  \textbf{\#} & \textbf{Declaración} & \textbf{Sí} \\
  \midrule
  A & Actividad docente del grupo \TTODfull & $\square$ \\
  B & Uso investigador secundario seudonimizado de artefactos de proceso & $\square$ \\
  C & Corpus investigador: sin categorías especiales; sin identificadores publicados & $\square$ \\
  D & Encuesta/entrevista (solo si se activa en el futuro) & N/A \\
  E & Grabación audio/vídeo de defensa o entrevista (solo si se activa) & N/A \\
  F & Difusión de extractos no identificativos (opcional) & $\square$ \\
  G & Imagen en LinkedIn/medios/publicaciones/revistas internas UDIT/web público (opcional) & $\square$ \\
  \bottomrule
\end{tabularx}

\vspace{0.8em}
\noindent\textbf{Nombre y apellidos:} \hrulefill

\vspace{0.7em}
\noindent
\begin{tabularx}{\textwidth}{@{}X X@{}}
  \textbf{Firma:} \hrulefill &
  \textbf{Fecha:} \underline{\hspace{1.1cm}}~/~\underline{\hspace{1.1cm}}~/~\underline{\hspace{2cm}}
\end{tabularx}

\vspace{0.6em}
\noindent{\scriptsize
Uso interno: seudónimo \texttt{P\_\_\_\_} · custodio: \underline{\hspace{3cm}} ·
archivo fuera de git.}

%======================================================================
% 7. EXCLUSIÓN FILOSÓFICA
%======================================================================
\AnnexTitle{7.\ Exclusión de la línea filosófica paralela}

Existe en el proyecto \TTODfull\ una línea de investigación filosófica
(ontología / epistemología / cibernética). Está
\textbf{aparada} y \textbf{no forma parte} de esta solicitud. \textbf{No usa}
datos de estudiantes ni artefactos de la Cohorte Front-End~II.

Esta remisión se limita al estudio pedagógico de caso descrito en los Anexos~1--6.

\vspace{1.2em}
\noindent{\color{rulegray}\hrule height 0.4pt}
\vspace{0.5em}
\AuthoritiesBlock

\vspace{0.6em}
\noindent{\footnotesize\color{uditnavy}
Fin del paquete de anexos. Formulario CEI Word:
\texttt{FORMULARIO-SOLICITUD-CEI-UDIT-RELLENO.docx}.}

\end{document}
